\ifdefined\gkver\else\def\gkver{student}\fi
\documentclass[\gkver]{gaokaozhenti}
\begin{document}

\chapter{2013年上海}

\begin{enumerate}
\item
%% number: 1
%% paperName: 2013年上海高考第 1 题 2 分
%% typeId: 1
%% type: 单选题
%% chapter: 光学
%% point: 光的电磁说
%% method: 无
%% score: 2
%% degree: 850
%% duplicateId: 0
%% body:
电磁波与机械波具有的共同性质是\xzanswer{B}
\fourchoices[answer=B]
{都是横波}
{都能传输能量}
{都能在真空中传播}
{都具有恒定的波速}
%% answer: B
%% memo:
\memoanswer{
机械波既有横波也有纵波，故 A 错误；电磁波和机械波都能传输能量，故 B 正确；机械波不能在真空中传播，故 C 错误；电磁波在真空中具有相同速度，但不同频率的电磁波在同一均匀介质中传播速度不同，而机械波在同一均匀介质中传播速度相同，故 D 错误。
}

\item
%% number: 2
%% paperName: 2013年上海高考第 2 题 2 分
%% typeId: 1
%% type: 单选题
%% chapter: 光学
%% point: 光电效应
%% method: 无
%% score: 2
%% degree: 850
%% duplicateId: 0
%% body:
当用一束紫外线照射锌板时，产生了光电效应，这时\xzanswer{C}
\fourchoices[answer=C]
{锌板带负电}
{有正离子从锌板逸出}
{有电子从锌板逸出}
{锌板会吸附空气中的正离子}
%% answer: C
%% memo:
\memoanswer{
发生光电效应时，有光电子从锌板中逸出，逸出光电子后的锌板带正电，对空气中的正离子有排斥作用，故 C 正确。
}

\item
%% number: 3
%% paperName: 2013年上海高考第 3 题 2 分
%% typeId: 1
%% type: 单选题
%% chapter: 光学
%% point: 光的干涉
%% method: 无
%% score: 2
%% degree: 850
%% duplicateId: 0
%% body:
白光通过双缝后产生的干涉条纹是彩色的，其原因是不同色光的\xzanswer{D}
\fourchoices[answer=D]
{传播速度不同}
{强度不同}
{振动方向不同}
{频率不同}
%% answer: D
%% memo:
\memoanswer{
光的颜色由光的频率决定，不同色光的波长不相同，故频率也不相同，光屏上不同位置出现不同色光，是该色光在此位置干涉加强的结果，故 D 正确。
}

\item
%% number: 4
%% paperName: 2013年上海高考第 4 题 2 分
%% typeId: 1
%% type: 单选题
%% chapter: 机械振动
%% point: 简谐振动
%% method: 无
%% score: 2
%% degree: 850
%% duplicateId: 0
%% body:
做简谐振动的物体，当它每次经过同一位置时，可能不同的物理量是\xzanswer{B}
\fourchoices[answer=B]
{位移}
{速度}
{加速度}
{回复力}
%% answer: B
%% memo:
\memoanswer{
做简谐振动的物体，经过同一位置时，相对平衡位置的位移 $x$ 相同，回复力（$F=-kx$）相同，由牛顿第二定律（$F=ma$）知加速度 $a$ 相同，物体可能以方向相反的速度通过同一点，故 B 正确。
}

\item
%% number: 5
%% paperName: 2013年上海高考第 5 题 2 分
%% typeId: 1
%% type: 单选题
%% chapter: 固体液体的性质
%% point: 固体的基本性质
%% method: 无
%% score: 2
%% degree: 850
%% duplicateId: 0
%% body:
液体与固体具有的相同特点是\xzanswer{B}
\fourchoices[answer=B]
{都具有确定的形状}
{体积都不易被压缩}
{物质分子的位置都确定}
{物质分子都在固定位置附近振动}
%% answer: B
%% memo:
\memoanswer{
液体没有确定的形状，故 A 错误；液体和固体分子间的距离都比较小，不易被压缩，故 B 正确；分子不停地做无规则的热运动，液体分子没有固定的位置，故 C、D 错误。
}

\item
%% number: 6
%% paperName: 2013年上海高考第 6 题 2 分
%% typeId: 1
%% type: 单选题
%% chapter: 曲线运动
%% point: 竖直平面圆周运动
%% method: 无
%% score: 2
%% degree: 650
%% duplicateId: 0
%% body:
秋千的吊绳有些磨损。在摆动过程中，吊绳最容易断裂的时候是秋千\xzanswer{D}
\fourchoices[answer=D]
{在下摆过程中}
{在上摆过程中}
{摆到最高点时}
{摆到最低点时}
%% answer: D
%% memo:
\memoanswer{
秋千在最低点时，速度最大，所需向心力最大，这时对绳的拉力也最大，故 D 正确。
}

\item
%% number: 7
%% paperName: 2013年上海高考第 7 题 2 分
%% typeId: 1
%% type: 单选题
%% chapter: 原子和原子核
%% point: 天然放射性
%% method: 无
%% score: 2
%% degree: 850
%% duplicateId: 0
%% body:
在一个 $\ce{^{238}_{92}U}$ 原子核衰变为一个 $\ce{^{206}_{82}Pb}$ 原子核的过程中，发生 $\beta$ 衰变的次数为\xzanswer{A}
\fourchoices[answer=A]
{6 次}
{10 次}
{22 次}
{32 次}
%% answer: A
%% memo:
\memoanswer{
发生 $\alpha$ 衰变的次数为 $\frac{238-206}{4}$ 次 $=8$ 次；发生 $\beta$ 衰变的次数为 $82-(92-2\times8)=6$ 次，故 A 正确。
}

\item
%% number: 8
%% paperName: 2013年上海高考第 8 题 2 分
%% typeId: 1
%% type: 单选题
%% chapter: 运动和力
%% point: 牛顿第二定律
%% method: 整体与隔离
%% score: 2
%% degree: 650
%% duplicateId: 0
%% body:
如图，质量 $m_{A} > m_{B}$ 的两物体 A、B 叠放在一起，靠着竖直墙面。让它们由静止释放，在沿粗糙墙面下落过程中，物体 B 的受力示意图是\xzanswer{A}
\onepicture[align=right, h=4.2cm]{\includesvg{figs/08}}
\fourchoices[answer=A, ispicture=true, h=2.2cm]
{figs/08a.png}
{figs/08b.png}
{figs/08c.png}
{figs/08d.png}
%% answer: A
%% memo:
\memoanswer{
由水平方向的受力平衡知，墙对 A、B 无弹力作用，则无摩擦力作用，故 C、D 错误；A、B 均做自由落体运动，它们之间无作用力，故 A 正确，B 错误。
}

\item
%% number: 9
%% paperName: 2013年上海高考第 9 题 3 分
%% typeId: 1
%% type: 单选题
%% chapter: 曲线运动
%% point: 天体运动
%% method: 无
%% score: 3
%% degree: 650
%% duplicateId: 0
%% body:
小行星绕恒星运动，恒星均匀地向四周辐射能量，质量缓慢减小，可认为小行星在绕恒星运动一周的过程中近似做圆周运动。则经过足够长的时间后，小行星运动的\xzanswer{A}
\fourchoices[answer=A]
{半径变大}
{速率变大}
{角速度变大}
{加速度变大}
%% answer: A
%% memo:
\memoanswer{
恒星质量减小，则对行星的引力减小，行星逐渐远离恒星，半径变大，故 A 正确；由 $G\frac{Mm}{r^{2}}=ma=m\frac{v^{2}}{r}=mr\omega^{2}$，解得 $a=\frac{GM}{r^{2}}$、$v=\sqrt{\frac{GM}{r}}$、$\omega=\sqrt{\frac{GM}{r^{3}}}$，可知恒星质量减小时，小行星的速率、角速度和加速度均减小，故 B、C、D 错误。
}

\item
%% number: 10
%% paperName: 2013年上海高考第 10 题 3 分
%% typeId: 1
%% type: 单选题
%% chapter: 电场
%% point: 等势面
%% method: 无
%% score: 3
%% degree: 650
%% duplicateId: 0
%% body:
两异种点电荷电场中的部分等势面如图所示，已知 A 点电势高于 B 点电势。若位于 a、b 处点电荷的电荷量大小分别为 $q_{a}$ 和 $q_{b}$，则\xzanswer{B}
\onepicture[align=right, w=4.5cm]{figs/10.png}
\fourchoices[answer=B]
{a 处为正电荷，$q_{a} < q_{b}$}
{a 处为正电荷，$q_{a} > q_{b}$}
{a 处为负电荷，$q_{a} < q_{b}$}
{a 处为负电荷，$q_{a} > q_{b}$}
%% answer: B
%% memo:
\memoanswer{
由题中等势面分布图及 A 点电势高于 B 点电势可知，a 处为正电荷，故 C、D 错误；等势面密的地方场强大，由点电荷场强公式 $E=\frac{kq}{r^{2}}$ 知 $q_{a}>q_{b}$，故 A 错误，B 正确。
}

\item
%% number: 11
%% paperName: 2013年上海高考第 11 题 3 分
%% typeId: 1
%% type: 单选题
%% chapter: 电磁感应
%% point: 楞次定律推广
%% method: 无
%% score: 3
%% degree: 650
%% duplicateId: 0
%% body:
如图，通电导线 MN 与单匝矩形线圈 abcd 共面，位置靠近 ab 且相互绝缘。当 MN 中电流突然减小时，线圈所受安培力的合力方向\xzanswer{B}
\onepicture[align=right, w=4.5cm]{\includesvg{figs/11}}
\fourchoices[answer=B]
{向左}
{向右}
{垂直纸面向外}
{垂直纸面向里}
%% answer: B
%% memo:
\memoanswer{
当 MN 中电流突然减小时，穿过线圈中的合磁通量减小，线圈中产生感应电流，由楞次定律知产生的感应电流将阻碍磁通量的减小，即线圈所受的安培力的方向向右，故 B 正确。
}

\item
%% number: 12
%% paperName: 2013年上海高考第 12 题 3 分
%% typeId: 1
%% type: 单选题
%% chapter: 电路
%% point: 逻辑电路
%% method: 无
%% score: 3
%% degree: 650
%% duplicateId: 0
%% body:
在车门报警电路中，两个按钮开关分别装在汽车的两扇门上，只要有开关处于断开状态，报警灯就发光。能实现此功能的电路是\xzanswer{D}
\fourchoices[answer=D, ispicture=true, h=2.6cm]
{figs/12a.png}
{figs/12b.png}
{figs/12c.png}
{figs/12d.png}
%% answer: D
%% memo:
\memoanswer{
实现题目要求的是“或”门电路，而选项 A、B 是“与”门电路，故 A、B 错误；选项 C、D 是“或”门电路，但选项 C 中 $\mathrm{S}_{1}$、$\mathrm{S}_{2}$ 均断开时，灯泡不亮，故能实现题目要求的只有 D，故 D 正确。
}

\item
%% number: 13
%% paperName: 2013年上海高考第 13 题 3 分
%% typeId: 1
%% type: 单选题
%% chapter: 磁场
%% point: 电流的磁场
%% method: 无
%% score: 3
%% degree: 650
%% duplicateId: 0
%% body:
如图，足够长的直线 ab 靠近通电螺线管，与螺线管平行。用磁传感器测量 ab 上各点的磁感应强度 $B$，在计算机屏幕上显示的大致图像是\xzanswer{C}
\onepicture[align=right, w=5.2cm]{\includesvg{figs/13}}
\fourchoices[answer=C, ispicture=true, h=2.2cm]
{figs/13a.png}
{figs/13b.png}
{figs/13c.png}
{figs/13d.png}
%% answer: C
%% memo:
\memoanswer{
通电螺线管的磁场分布与条形磁铁的磁场分布类似，两磁极磁感应强度最大，远离磁极，磁感应强度减小，故 C 正确。
}

\item
%% number: 14
%% paperName: 2013年上海高考第 14 题 3 分
%% typeId: 1
%% type: 单选题
%% chapter: 机械波
%% point: 波长频率波速
%% method: 无
%% score: 3
%% degree: 650
%% duplicateId: 0
%% body:
一列横波沿水平绳传播，绳的一端在 $t = 0$ 时开始做周期为 $T$ 的简谐运动，经过时间 $t$（$\frac{3}{4}T < t < T$），绳上某点位于平衡位置上方的最大位移处。则在 $2t$ 时，该点位于平衡位置的\xzanswer{A}
\fourchoices[answer=A]
{上方，且向上运动}
{上方，且向下运动}
{下方，且向上运动}
{下方，且向下运动}
%% answer: A
%% memo:
\memoanswer{
位于平衡位置上方最大位移处的质点，再经 $\frac{3}{4}T\sim T$ 的时间，质点运动到平衡位置的上方，且向上运动，故 A 正确。
}

\item
%% number: 15
%% paperName: 2013年上海高考第 15 题 3 分
%% typeId: 1
%% type: 单选题
%% chapter: 气体性质
%% point: 理想气体状态方程
%% method: 无
%% score: 3
%% degree: 650
%% duplicateId: 0
%% body:
已知湖水深度为 $20\Um$，湖底水温为 $4\UduC$，水面温度为 $17\UduC$，大气压强为 $1.0\times10^{5}\UPa$。当一气泡从湖底缓慢升到水面时，其体积约为原来的（取 $g = 10\Umsq$，$\rho_{\text{水}} = 1.0\times10^{3}\Ukgmc$）\xzanswer{C}
\fourchoices[answer=C]
{12.8 倍}
{8.5 倍}
{3.1 倍}
{2.1 倍}
%% answer: C
%% memo:
\memoanswer{
由理想气体状态方程 $\frac{p_{1}V_{1}}{T_{1}}=\frac{p_{2}V_{2}}{T_{2}}$ 解得 $\frac{V_{2}}{V_{1}}=\frac{p_{1}T_{2}}{p_{2}T_{1}}=\frac{(1.0\times10^{5}+1.0\times10^{3}\times10\times20)\times290}{1.0\times10^{5}\times277}\approx 3.1$，故 C 正确。
}

\item
%% number: 16
%% paperName: 2013年上海高考第 16 题 3 分
%% typeId: 1
%% type: 单选题
%% chapter: 功和能
%% point: 交通工具启动
%% method: 无
%% score: 3
%% degree: 650
%% duplicateId: 0
%% body:
汽车以恒定功率沿公路做直线运动，途中通过一块沙地。汽车在公路及沙地上所受阻力均为恒力，且在沙地上受到的阻力大于在公路上受到的阻力。汽车在驶入沙地前已做匀速直线运动，它在驶入沙地到驶出沙地后的一段时间内，位移 $s$ 随时间 $t$ 的变化关系可能是\xzanswer{A}
\fourchoices[answer=A, ispicture=true, h=2.6cm]
{figs/16a.png}
{figs/16b.png}
{figs/16c.png}
{figs/16d.png}
%% answer: A
%% memo:
\memoanswer{
汽车在恒定功率下运动，驶入沙地前有 $F_{\text{牵}}=f$；驶入沙地后，$F_{\text{牵}}<f_{\text{沙}}$，汽车做减速运动，又 $P=F_{\text{牵}}v$，$v$ 减小，$F_{\text{牵}}$ 增大，当 $F_{\text{牵}}=f_{\text{沙}}$ 时，汽车做匀速运动；驶出沙地后有 $F_{\text{牵}}>f$，汽车做加速运动。$s-t$ 图象中图线的斜率表示速度的大小，故 A 正确。
}

\item
%% number: 17
%% paperName: 2013年上海高考第 17 题 4 分
%% typeId: 2
%% type: 多选题
%% chapter: 光学
%% point: 光子说
%% method: 无
%% score: 4
%% degree: 650
%% duplicateId: 0
%% body:
某半导体激光器发射波长为 $1.5\times10^{-6}\Um$，功率为 $5.0\times10^{-3}\UW$ 的连续激光。已知可见光波长的数量级为 $10^{-7}\Um$，普朗克常量 $h = 6.63\times10^{-34}\UJcdots$，该激光器发出的\xzanswer{BD}
\fourchoices[answer=BD]
{是紫外线}
{是红外线}
{光子能量约为 $1.3\times10^{-18}\UJ$}
{光子数约为每秒 $3.8\times10^{16}$ 个}
%% answer: BD
%% memo:
\memoanswer{
由于激光器发出的光波波长比可见光长，所以发出的是红外线，故 A 错误，B 正确；光子能量为 $E = h\nu = h\frac{c}{\lambda} = 1.3 \times 10^{-19}\UJ$，故 C 错误；每秒发射的光子数约为 $n = \frac{P\cdot 1\Us}{E} = 3.8 \times 10^{16}$ 个，故 D 正确。
}

\item
%% number: 18
%% paperName: 2013年上海高考第 18 题 4 分
%% typeId: 2
%% type: 多选题
%% chapter: 力物体的平衡
%% point: 力的合成
%% method: 无
%% score: 4
%% degree: 650
%% duplicateId: 0
%% body:
两个共点力 $F_{1}$、$F_{2}$ 大小不同，它们的合力大小为 $F$，则\xzanswer{AD}
\fourchoices[answer=AD]
{$F_{1}$、$F_{2}$ 同时增大一倍，$F$ 也增大一倍}
{$F_{1}$、$F_{2}$ 同时增加 $10\UN$，$F$ 也增加 $10\UN$}
{$F_{1}$ 增加 $10\UN$，$F_{2}$ 减少 $10\UN$，$F$ 一定不变}
{若 $F_{1}$、$F_{2}$ 中的一个增大，$F$ 不一定增大}
%% answer: AD
%% memo:
\memoanswer{
通过作平行四边形可知 A 正确；通过分析同一直线的两个力 $F_{1}$、$F_{2}$ 的合成可知，故 B、C 错误，D 正确。
}

\item
%% number: 19
%% paperName: 2013年上海高考第 19 题 4 分
%% typeId: 2
%% type: 多选题
%% chapter: 曲线运动
%% point: 平抛运动
%% method: 无
%% score: 4
%% degree: 450
%% duplicateId: 0
%% body:
如图，轰炸机沿水平方向匀速飞行，到达山坡底端正上方时释放一颗炸弹，并垂直击中山坡上的目标 A。已知 A 点高度为 $h$，山坡倾角为 $\theta$，由此可算出\xzanswer{ABC}
\onepicture[align=right, w=4.5cm]{\includesvg{figs/19}}
\fourchoices[answer=ABC]
{轰炸机的飞行高度}
{轰炸机的飞行速度}
{炸弹的飞行时间}
{炸弹投出时的动能}
%% answer: ABC
%% memo:
\memoanswer{
炸弹做平抛运动，竖直分位移 $y = \frac{1}{2}v_{y}t = \frac{1}{2}\cdot\frac{v_{0}}{\tan\theta}t = \frac{1}{2\tan\theta}\cdot v_{0}t = \frac{x}{2\tan\theta} = \frac{h}{2\tan^{2}\theta}$，则轰炸机的飞行高度为 $H = y + h = h\left(\frac{1}{2\tan^{2}\theta} + 1\right)$，故 A 正确；由 $y = \frac{1}{2}gt^{2}$ 可求炸弹的飞行时间 $t$，故 C 正确；由水平位移 $x = \frac{h}{\tan\theta} = v_{0}t$ 可求炸弹的水平初速度，即轰炸机的飞行速度，故 B 正确；由于不知炸弹的质量，故无法确定炸弹投出时的动能，故 D 错误。
}

\item
%% number: 20
%% paperName: 2013年上海高考第 20 题 4 分
%% typeId: 2
%% type: 多选题
%% chapter: 直线运动
%% point: 运动的分解
%% method: 无
%% score: 4
%% degree: 450
%% duplicateId: 0
%% body:
右图为在平静海面上，两艘拖船 A、B 拖着驳船 C 运动的示意图。A、B 的速度分别沿着缆绳 CA、CB 方向，A、B、C 不在一条直线上。由于缆绳不可伸长，因此 C 的速度在 CA、CB 方向的投影分别与 A、B 的速度相等，由此可知 C 的\xzanswer{BC}
\onepicture[align=right, h=4.5cm]{\includesvg{figs/20}}
\fourchoices[answer=BC]
{速度大小可以介于 A、B 的速度大小之间}
{速度大小一定不小于 A、B 的速度大小}
{速度方向可能在 CA 和 CB 的夹角范围外}
{速度方向一定在 CA 和 CB 的夹角范围内}
%% answer: BC
%% memo:
\memoanswer{
船 C 沿着绳子靠向 A 船初位置运动的同时还要绕 A 船转动，船 C 沿着绳子靠向 B 船的初位置运动的同时还要绕 B 船转动，先将船 C 的速度沿着平行于 AC 绳子和垂直于 AC 绳子方向正交分解，再将船 C 的速度沿着平行于 BC 绳子和垂直于 BC 绳子方向正交分解，如图所示。

\onepicture[w=3cm]{figs/20_an.jpg}

由于绳子不可伸长，故每条船沿着绳子方向的分速度是相等的；两拖船速度一定小于 C 船速度（斜边大于直角边），故 A 错误，B 正确；因 C 的速度在 CA、CB 方向的投影分别与 A、B 的速度相等，速度方向不一定在 CA 和 CB 的夹角范围内，以 A 为参照系，此时的 C 船是以 A 为圆心做圆周运动，即速度方向垂直于 AC，$\angle ACB$ 小于 $90^{\circ}$ 的情况下，C 的速度方向就是在 ACB 夹角之外的，故 C 正确，D 错误。综上所述 B、C 正确。
}

\item
%% number: 21
%% paperName: 2013年上海高考第 21 题 4 分
%% typeId: 3
%% type: 填空题
%% chapter: 原子和原子核
%% point: 天然放射性
%% method: 无
%% score: 4
%% degree: 850
%% duplicateId: 0
%% body:
放射性元素 $\ce{^{210}_{84}Po}$ 衰变为 $\ce{^{206}_{82}Pb}$，此衰变过程的核反应方程是\tkanswer[4.5cm]{$\ce{^{210}_{84}Po -> ^{206}_{82}Pb + ^{4}_{2}He}$}；用此衰变过程中发出的射线轰击 $\ce{^{19}_{9}F}$，可得到质量数为 22 的氖（Ne）元素和另一种粒子，此核反应过程的方程是\tkanswer[4.5cm]{$\ce{^{19}_{9}F + ^{4}_{2}He -> ^{22}_{10}Ne + ^{1}_{1}H}$}。
%% answer: $\ce{^{210}_{84}Po -> ^{206}_{82}Pb + ^{4}_{2}He}$，$\ce{^{19}_{9}F + ^{4}_{2}He -> ^{22}_{10}Ne + ^{1}_{1}H}$
%% memo:
\memoanswer{
由电荷数守恒和质量数守恒得 $\ce{^{210}_{84}Po -> ^{206}_{82}Pb + ^{4}_{2}He}$，$\ce{^{19}_{9}F + ^{4}_{2}He -> ^{22}_{10}Ne + ^{1}_{1}H}$。
}

\item
%% number: 22
%% paperName: 2013年上海高考第 22 题 4 分
%% typeId: 3
%% type: 填空题
%% chapter: 曲线运动
%% point: 人造地球卫星
%% method: 无
%% score: 4
%% degree: 650
%% duplicateId: 0
%% body:
（本大题中第 22 题为分叉题，分 A、B 两类，考生可任选一类答题。若两类试题均做，一律按 A 类题计分。）

22A．质量为 $M$ 的物块静止在光滑水平桌面上，质量为 $m$ 的子弹以水平速度 $v_{0}$ 射入物块后，以水平速度 $\frac{2}{3}v_{0}$ 射出。则物块的速度为\tkanswer[2cm]{$\frac{m}{3M}v_{0}$}，此过程中损失的机械能为\tkanswer[3.5cm]{$\left(\frac{5}{18}-\frac{m}{18M}\right)mv_{0}^{2}$}。

22B．若两颗人造地球卫星的周期之比为 $T_{1}:T_{2} = 2:1$，则它们的轨道半径之比 $R_{1}:R_{2} =$\tkanswer[2cm]{$\sqrt[3]{4}:1$}，向心加速度之比 $a_{1}:a_{2} =$\tkanswer[2.5cm]{$1:\sqrt[3]{16}$}。
%% answer: $\frac{m}{3M}v_{0}$，$\left(\frac{5}{18}-\frac{m}{18M}\right)mv_{0}^{2}$；$\sqrt[3]{4}:1$，$1:\sqrt[3]{16}$
%% memo:
\memoanswer{
22A．由动量守恒得 $mv_{0}=Mv+m\cdot\frac{2}{3}v_{0}$，解得 $v=\frac{mv_{0}}{3M}$；由能量守恒知损失的机械能为 $\Delta E=\frac{1}{2}mv_{0}^{2}-\frac{1}{2}Mv^{2}-\frac{1}{2}m\cdot\left(\frac{2}{3}v_{0}\right)^{2}=\left(\frac{5}{18}-\frac{m}{18M}\right)mv_{0}^{2}$。

22B．由 $G\frac{Mm}{R^{2}}=mR\left(\frac{2\pi}{T}\right)^{2}$ 解得 $R=\sqrt[3]{\frac{GMT^{2}}{4\pi^{2}}}$，则 $R_{1}:R_{2}=\sqrt[3]{T_{1}^{2}}:\sqrt[3]{T_{2}^{2}}=\sqrt[3]{4}:1$；由 $a=R\left(\frac{2\pi}{T}\right)^{2}$ 得 $a_{1}:a_{2}=R_{1}T_{2}^{2}:R_{2}T_{1}^{2}=1:2\sqrt[3]{2}$。
}

\item
%% number: 23
%% paperName: 2013年上海高考第 23 题 4 分
%% typeId: 3
%% type: 填空题
%% chapter: 机械振动
%% point: 等效单摆
%% method: 无
%% score: 4
%% degree: 650
%% duplicateId: 0
%% body:
如图，在半径为 $2.5\Um$ 的光滑圆环上切下一小段圆弧，放置于竖直平面内，两端点距最低点高度差 $H$ 为 $1\Ucm$。将小环置于圆弧端点并从静止释放，小环运动到最低点所需的最短时间为\tkanswer[1.5cm]{$\frac{\pi}{4} \approx 0.785$}$\Us$，在最低点处的加速度为\tkanswer[1.2cm]{0.08}$\Umsq$。（取 $g = 10\Umsq$）
\onepicture[align=right, w=5cm]{\includesvg{figs/23}}
%% answer: $\frac{\pi}{4} \approx 0.785$，0.08
%% memo:
\memoanswer{
设圆弧半径为 $R$，由于 $H \ll R$，小环的运动可视为单摆摆动，则小环运动到最低点所需的最短时间为 $t = \frac{1}{4}T = \frac{1}{2}\pi\sqrt{\frac{R}{g}} = 0.785\Us$；由 $mgH = \frac{1}{2}mv^{2}$ 和 $a = \frac{v^{2}}{R}$ 联立解得 $a = \frac{2gH}{R} = 0.08\Umsq$。
}

\item
%% number: 24
%% paperName: 2013年上海高考第 24 题 4 分
%% typeId: 3
%% type: 填空题
%% chapter: 电路
%% point: 闭合电路欧姆定律
%% method: 无
%% score: 4
%% degree: 650
%% duplicateId: 0
%% body:
如图，电路中三个电阻 $R_{1}$、$R_{2}$ 和 $R_{3}$ 的阻值分别为 $R$、$2R$ 和 $4R$。当电键 $\mathrm{S}_{1}$ 断开、$\mathrm{S}_{2}$ 闭合时，电源输出功率为 $P_{0}$；当 $\mathrm{S}_{1}$ 闭合、$\mathrm{S}_{2}$ 断开时，电源输出功率也为 $P_{0}$。则电源电动势为\tkanswer[2.5cm]{$3\sqrt{P_{0}R}$}；当 $\mathrm{S}_{1}$、$\mathrm{S}_{2}$ 都断开时，电源的总功率为\tkanswer[1.5cm]{$P_{0}$}。
\onepicture[align=right, w=4.5cm]{\includesvg{figs/24}}
%% answer: $3\sqrt{P_{0}R}$，$P_{0}$
%% memo:
\memoanswer{
（1）当电键 $\mathrm{S}_{1}$ 断开、$\mathrm{S}_{2}$ 闭合时，$R_{2}$、$R_{3}$ 被短路，外电路只有 $R_{1}$ 一个电阻，有：$P_{0}=\left(\frac{E}{R_{1}+r}\right)^{2}R_{1}=\left(\frac{E}{R+r}\right)^{2}R$。当电键 $\mathrm{S}_{1}$ 闭合、$\mathrm{S}_{2}$ 断开时，$R_{1}$、$R_{2}$ 被短路，外电路只有 $R_{3}$ 一个电阻，有：$P_{0}=\left(\frac{E}{R_{3}+r}\right)^{2}R_{3}=\left(\frac{E}{4R+r}\right)^{2}4R$。联立以上两式，可求得 $E=3\sqrt{P_{0}R}$，$r=2R$。

（2）当 $\mathrm{S}_{1}$、$\mathrm{S}_{2}$ 都断开时，外电阻为 $R_{1}$、$R_{2}$、$R_{3}$ 三者串联，有：$P_{\text{总}}=\frac{E^{2}}{R_{1}+R_{2}+R_{3}+r}=\frac{(3\sqrt{P_{0}R})^{2}}{9R}=P_{0}$。
}

\item
%% number: 25
%% paperName: 2013年上海高考第 25 题 4 分
%% typeId: 3
%% type: 填空题
%% chapter: 力物体的平衡
%% point: 平衡综合
%% method: 无
%% score: 4
%% degree: 450
%% duplicateId: 0
%% body:
如图，倾角为 $37^\circ$，质量不计的支架 ABCD 的 D 端有一大小与质量均可忽略的光滑定滑轮，A 点处有一固定转轴，$\mathrm{CA}\perp\mathrm{AB}$，$\mathrm{DC}=\mathrm{CA}=0.3\Um$。质量 $m = 1\Ukg$ 的物体置于支架的 B 端，并与跨过定滑轮的轻绳相连，绳另一端作用一竖直向下的拉力 $F$，物体在拉力作用下沿 BD 做匀速直线运动，已知物体与 BD 间的动摩擦因数 $\mu = 0.3$。为保证支架不绕 A 点转动，物体向上滑行的最大距离 $s =$\tkanswer[1.5cm]{0.248}$\Um$。若增大 $F$ 后，支架仍不绕 A 点转动，物体能向上滑行的最大距离 $s'$\tkanswer[1cm]{等于}$s$（填：“大于”、“等于”或“小于”。）（取 $\sin37^\circ = 0.6$，$\cos37^\circ = 0.8$）
\onepicture[align=right, w=5cm]{\includesvg{figs/25}}
%% answer: 0.248，等于
%% memo:
\memoanswer{
以 $m$ 为研究对象，$N = mg\cos37^\circ = 8\UN$，$f = \mu N = 2.4\UN$，$F = mg\sin37^\circ + f = 8.4\UN$。以支架为研究对象，以 A 为支点，由几何关系知力 $F$ 的力臂为 $l_{1} = 0.24\Um$，摩擦力和平行斜面细线拉力的力臂为 $l_{2} = 0.24\Um$，压力的力臂为 $l_{3} = 0.32\Um - s$，由力矩平衡得 $Fl_{1} + fl_{2} = Nl_{3} + Fl_{2}$，解得 $s = 0.248\Um$。由于 $l_{1} = l_{2}$，可见增大 $F$ 后，物体向上滑行的最大距离不变，即 $s' = s$。
}

\item
%% number: 26
%% paperName: 2013年上海高考第 26 题 3 分
%% typeId: 4
%% type: 实验题
%% chapter: 电磁感应
%% point: 切割磁感线电动势
%% method: 无
%% score: 3
%% degree: 850
%% duplicateId: 0
%% body:
演示地磁场存在的实验装置（由环形线圈，微电流传感器，DIS 等组成）如图所示。首先将线圈竖直放置，以竖直方向为轴转动，屏幕上的电流指针\tkanswer[1cm]{有}（填：“有”或“无”）偏转；然后仍将线圈竖直放置，使其平面与东西向平行，并从东向西移动，电流指针\tkanswer[1cm]{无}（填：“有”或“无”）偏转；最后将线圈水平放置，使其从东向西移动，电流指针\tkanswer[1cm]{无}（填：“有”或“无”）偏转。
\onepicture[align=right, w=4.5cm]{figs/26.jpg}
%% answer: 有，无，无
%% memo:
\memoanswer{
地磁场方向为南北方向，闭合回路中的磁通量发生变化，则回路中产生感应电流；第一种情况线圈中的磁通量发生变化，后两种情况线圈中的磁通量没有发生变化。
}

\item
%% number: 27
%% paperName: 2013年上海高考第 27 题 6 分
%% typeId: 4
%% type: 实验题
%% chapter: 电路
%% point: 多用表的使用
%% method: 无
%% score: 6
%% degree: 650
%% duplicateId: 0
%% body:
为确定某电子元件的电气特性，做如下测量。
\begin{enumerate}
	\item
	用多用表测量该元件的电阻，选用“$\times$100”倍率的电阻档测量，发现多用表指针偏转过大，因此需选择\tkanswer[1.2cm]{$\times$10}倍率的电阻档（填：“$\times$10”或“$\times$1k”），并\tkanswer[1.5cm]{重新调零}再进行测量，多用表的示数如图（a）所示，测量结果为\tkanswer[1.5cm]{70}$\UO$。

	\item
	将待测元件（额定电压 $9\UV$）、蓄电池、滑动变阻器、电流表、多用表、电键及若干导线连接成电路如图（b）所示。添加连线，使电路能测量该元件完整的伏安特性。本实验中使用多用表测电压，多用表的选择开关应调到\tkanswer[2.6cm]{直流电压 $10\UV$}档（填：“直流电压 $10\UV$”或“直流电压 $50\UV$”）。
\end{enumerate}
\twopicture[align=right, wA=5cm, wB=4.2cm]{figs/27a.png}{\includesvg{figs/27b}}
%% answer: $\times$10，重新调零，70$\UO$；连线如图，直流电压 $10\UV$
\jdanswer{
\begin{enumerate}
	\item
	$\times$10，重新调零，$70\UO$

	\item
	连线如图，直流电压 $10\UV$

\end{enumerate}
}
%% memo:
\memoanswer{
（1）用多用表“$\times$100”倍率测电阻时，指针偏转角度过大，示数太小，应换小倍率，即选择“$\times$10”倍率，欧姆挡调零后重新测量，测量结果为 $7\times10\UO=70\UO$。

（2）实物连接如答案图所示。由于待测元件额定电压为 $9\UV$，因此多用表的选择开关应调到直流电压 $10\UV$ 挡。

\onepicture[w=5cm]{\includesvg{figs/27_an}}
}

\item
%% number: 28
%% paperName: 2013年上海高考第 28 题 8 分
%% typeId: 4
%% type: 实验题
%% chapter: 曲线运动
%% point: 研究平抛运动实验
%% method: 无
%% score: 8
%% degree: 650
%% duplicateId: 0
%% body:
如图，研究平抛运动规律的实验装置放置在水平桌面上，利用光电门传感器和碰撞传感器可测得小球的水平初速度和飞行时间，底板上的标尺可以测得水平位移。保持水平槽口距底板高度 $h = 0.420\Um$ 不变。改变小球在斜槽导轨上下滑的起始位置，测出小球做平抛运动的初速度 $v_{0}$、飞行时间 $t$ 和水平位移 $d$，记录在表中。

\begin{center}
\begin{tabular}{|c|c|c|c|c|}
\hline
$v_{0}$（$\Ums$） & 0.741 & 1.034 & 1.318 & 1.584 \\ \hline
$t$（$\UmsT$） & 292.7 & 293.0 & 292.8 & 292.9 \\ \hline
$d$（$\Ucm$） & 21.7 & 30.3 & 38.6 & 46.4 \\ \hline
\end{tabular}
\end{center}

\begin{enumerate}
	\item
	由表中数据可知，在 $h$ 一定时，小球水平位移 $d$ 与其初速度 $v_{0}$ 成\tkanswer[1cm]{正比}关系，与\tkanswer[1.3cm]{飞行时间 $t$}无关。

	\item
	一位同学计算出小球飞行时间的理论值 $t_{\text{理}} = \sqrt{\frac{2h}{g}} = \sqrt{\frac{2 \times 0.420}{10}} = 289.8\UmsT$，发现理论值与测量值之差约为 $3\UmsT$。经检查，实验及测量无误，其原因是\tkanswer[3cm]{重力加速度值取值不正确}。

	\item
	另一位同学分析并纠正了上述偏差后，另做了这个实验，竟发现测量值 $t'$ 依然大于自己得到的理论值 $t_{\text{理}}'$，但二者之差在 $3\sim7\UmsT$ 之间，且初速度越大差值越小。对实验装置的安装进行检查，确认斜槽槽口与底座均水平，则导致偏差的原因是\tkanswer[4.5cm]{光电门传感器置于槽口的内侧}。
\end{enumerate}
\onepicture[align=right, w=6cm]{\includesvg{figs/28}}
%% answer: 正比，飞行时间 $t$；重力加速度值取值不正确；光电门传感器置于槽口的内侧
\jdanswer{
\begin{enumerate}
	\item
	正比，飞行时间 $t$

	\item
	重力加速度值取值不正确

	\item
	光电门传感器置于槽口的内侧（或小球直径过大、小球飞过光电门需要时间）

\end{enumerate}
}
%% memo:
\memoanswer{
（1）由表中数据可知小球水平位移 $d$ 与其初速度 $v_{0}$ 成正比，与 $h$ 无关。

（2）由题中计算小球飞出时间时选取 $g=10\Umsq$ 可知，实验误差可能是由于所选取的重力加速度不正确引起的。

（3）光电门传感器置于槽口的内侧（或小球直径过大、小球飞过光电门需要时间）。
}

\item
%% number: 29
%% paperName: 2013年上海高考第 29 题 7 分
%% typeId: 4
%% type: 实验题
%% chapter: 气体性质
%% point: 验证玻意耳定律
%% method: 无
%% score: 7
%% degree: 450
%% duplicateId: 0
%% body:
利用如图装置可测量大气压强和容器的容积。步骤如下：

\begin{steps}
	\item
	将倒 U 形玻璃管 A 的一端通过橡胶软管与直玻璃管 B 连接，并注入适量的水，另一端插入橡皮塞，然后塞住烧瓶口，并在 A 上标注此时水面的位置 K；再将一活塞置于 $10\UmL$ 位置的针筒插入烧瓶，使活塞缓慢推移至 0 刻度位置；上下移动 B，保持 A 中的水面位于 K 处，测得此时水面的高度差为 $17.1\Ucm$。

	\item
	拔出橡皮塞，将针筒活塞置于 $0\UmL$ 位置，使烧瓶与大气相通后再次塞住瓶口；然后将活塞抽拔至 $10\UmL$ 位置，上下移动 B，使 A 中的水面仍位于 K，测得此时玻璃管中水面的高度差为 $16.8\Ucm$。（玻璃管 A 内气体体积忽略不计，$\rho_{\text{水}} = 1.0\times10^{3}\Ukgmc$，取 $g = 10\Umsq$）
\end{steps}

\begin{enumerate}
	\item
	若用 $V_{0}$ 表示烧瓶容积，$p_{0}$ 表示大气压强，$\Delta V$ 表示针筒内气体的体积，$\Delta p_{1}$、$\Delta p_{2}$ 表示上述步骤①、②中烧瓶内外气体压强差大小，则步骤①、②中，气体满足的方程分别为\tkanswer[4.2cm]{$p_{0}(V_{0}+\Delta V)=(p_{0}+\Delta p_{1})V_{0}$}、\tkanswer[4.4cm]{$p_{0}V_{0}=(p_{0}-\Delta p_{2})(V_{0}+\Delta V)$}。

	\item
	由实验数据得烧瓶容积 $V_{0} =$\tkanswer[1.2cm]{560}$\UmL$，大气压强 $p_{0} =$\tkanswer[2cm]{$9.58\times10^{4}$}$\UPa$。

	\item
	（单选题）倒 U 形玻璃管 A 内气体的存在\xzanswer{A}
	\fourchoices[answer=A]
	{仅对容积的测量结果有影响}
	{仅对压强的测量结果有影响}
	{对二者的测量结果均有影响}
	{对二者的测量结果均无影响}
\end{enumerate}
\onepicture[align=right, w=4.2cm]{\includesvg{figs/29}}
%% answer: $p_{0}(V_{0}+\Delta V)=(p_{0}+\Delta p_{1})V_{0}$，$p_{0}V_{0}=(p_{0}-\Delta p_{2})(V_{0}+\Delta V)$；560，$9.58\times10^{4}$；A
\jdanswer{
\begin{enumerate}
	\item
	$p_{0}(V_{0}+\Delta V)=(p_{0}+\Delta p_{1})V_{0}$，$p_{0}V_{0}=(p_{0}-\Delta p_{2})(V_{0}+\Delta V)$

	\item
	$V_{0}=560\UmL$，$p_{0}=9.58\times10^{4}\UPa$

	\item
	A

\end{enumerate}
}
%% memo:
\memoanswer{
（1）步骤①为等温过程，满足方程 $p_{0}(V_{0}+\Delta V)=(p_{0}+\Delta p_{1})V_{0}$；步骤②也是等温过程，满足方程 $p_{0}V_{0}=(p_{0}-\Delta p_{2})(V_{0}+\Delta V)$。

（2）将已知数据代入以上两方程，联立解得 $V_{0}=560\UmL$，$p_{0}=9.58\times10^{4}\UPa$。

（3）倒 U 形玻璃管 A 内气体的存在仅对容积的测量结果有影响，故 A 正确。
}

\item
%% number: 30
%% paperName: 2013年上海高考第 30 题 10 分
%% typeId: 6
%% type: 计算题
%% chapter: 气体性质
%% point: 理想气体状态方程
%% method: 无
%% score: 10
%% degree: 650
%% duplicateId: 0
%% body:
如图，柱形容器内用不漏气的轻质绝热活塞封闭一定量的理想气体，容器外包裹保温材料。开始时活塞至容器底部的高度为 $H_{1}$，容器内气体温度与外界温度相等。在活塞上逐步加上多个砝码后，活塞下降到距容器底部 $H_{2}$ 处，气体温度升高了 $\Delta T$；然后取走容器外的保温材料，活塞位置继续下降，最后静止于距容器底部 $H_{3}$ 处：已知大气压强为 $p_{0}$。求：气体最后的压强与温度。
\onepicture[align=right, w=5cm]{\includesvg{figs/30}}
%% answer: $p_{3} = \frac{H_{1}}{H_{3}}p_{0}$，$T_{3} = \frac{H_{3}}{H_{2}-H_{3}}\Delta T$
\jdanswer{
$p_{3} = \frac{H_{1}}{H_{3}}p_{0}$，$T_{3} = \frac{H_{3}}{H_{2}-H_{3}}\Delta T$
}
%% memo:
\memoanswer{
解：以柱形容器内的气体为研究对象，由于气体的初末状态温度相等，由

$p_{0}H_{1}S = p_{3}H_{3}S$

得 $p_{3} = \frac{H_{1}}{H_{3}}p_{0}$

活塞位于 $H_{2}$ 时，温度 $T_{2} = T_{1} + \Delta T$，然后等压压缩至 $H_{3}$，且 $T_{3} = T_{1}$，由

$\frac{H_{2}S}{T_{2}} = \frac{H_{3}S}{T_{3}}$

得 $T_{3}H_{2} = (T_{3} + \Delta T)H_{3}$

解得 $T_{3} = \frac{H_{3}}{H_{2} - H_{3}}\Delta T$
}

\item
%% number: 31
%% paperName: 2013年上海高考第 31 题 12 分
%% typeId: 6
%% type: 计算题
%% chapter: 运动和力
%% point: 牛顿定律的应用
%% method: 分情况讨论
%% score: 12
%% degree: 450
%% duplicateId: 0
%% body:
如图，质量为 $M$、长为 $L$、高为 $h$ 的矩形滑块置于水平地面上，滑块与地面间动摩擦因数为 $\mu$；滑块上表面光滑，其右端放置一个质量为 $m$ 的小球。用水平外力击打滑块左端，使其在极短时间内获得向右的速度 $v_{0}$，经过一段时间后小球落地。求小球落地时距滑块左端的水平距离。
\onepicture[align=right, w=5cm]{\includesvg{figs/31}}
%% answer: 若 $\sqrt{\frac{2h}{g}} < \frac{1}{\mu g}\sqrt{v_{0}^{2} - 2\mu\left(1+\frac{m}{M}\right)gL}$，$s = \sqrt{\frac{2h}{g}\left[v_{0}^{2} - 2\mu\left(1+\frac{m}{M}\right)gL\right]} - \mu h$；若 $\sqrt{\frac{2h}{g}} > \frac{1}{\mu g}\sqrt{v_{0}^{2} - 2\mu\left(1+\frac{m}{M}\right)gL}$，$s' = \frac{v_{0}^{2}}{2\mu g} - \left(1+\frac{m}{M}\right)L$
\jdanswer{
若 $\sqrt{\frac{2h}{g}} < \frac{1}{\mu g}\sqrt{v_{0}^{2} - 2\mu\left(1+\frac{m}{M}\right)gL}$，$s = \sqrt{\frac{2h}{g}\left[v_{0}^{2} - 2\mu\left(1+\frac{m}{M}\right)gL\right]} - \mu h$；若 $\sqrt{\frac{2h}{g}} > \frac{1}{\mu g}\sqrt{v_{0}^{2} - 2\mu\left(1+\frac{m}{M}\right)gL}$，$s' = \frac{v_{0}^{2}}{2\mu g} - \left(1+\frac{m}{M}\right)L$。
}
%% memo:
\memoanswer{
解：小球在滑块上运动时，滑块运动满足 $\mu(M+m)g=Ma$，因此滑块的加速度大小

$a=\mu\left(1+\frac{m}{M}\right)g$

由 $v^{2}=v_{0}^{2}-2aL$ 得小球脱离滑块时滑块的速度

$v=\sqrt{v_{0}^{2}-2\mu\left(1+\frac{m}{M}\right)gL}$

小球脱离滑块后做自由落体运动，$t=\sqrt{\frac{2h}{g}}$

小球脱离滑块后滑块加速度 $a'=\mu g$

滑块的运动时间 $t'=\frac{v}{\mu g}=\frac{1}{\mu g}\sqrt{v_{0}^{2}-2\mu\left(1+\frac{m}{M}\right)gL}$

若 $t<t'$，小球落地时滑块尚未停止运动，在时间 $t$ 内滑块运动距离即小球落地时距滑块左侧的距离

$s=vt-\frac{1}{2}a't^{2}=\sqrt{\frac{2h}{g}\left[v_{0}^{2}-2\mu\left(1+\frac{m}{M}\right)gL\right]}-\mu h$

若 $t>t'$，小球落地前滑块已停止运动，则由 $0=v^{2}-2a's'$ 得

$v_{0}^{2}-2\mu\left(1+\frac{m}{M}\right)gL=2\mu gs'$

$s'=\frac{v_{0}^{2}}{2\mu g}-\left(1+\frac{m}{M}\right)L$
}

\item
%% number: 32
%% paperName: 2013年上海高考第 32 题 12 分
%% typeId: 6
%% type: 计算题
%% chapter: 电场
%% point: 带电粒子的直线运动
%% method: 无
%% score: 12
%% degree: 650
%% duplicateId: 0
%% body:
半径为 $R$，均匀带正电荷的球体在空间产生球对称的电场；场强大小沿半径分布如图所示，图中 $E_{0}$ 已知，$E-r$ 曲线下 $O-R$ 部分的面积等于 $R-2R$ 部分的面积。
\begin{enumerate}
	\item
	写出 $E-r$ 曲线下面积的单位；
	\item
	已知带电球在 $r \geq R$ 处的场强 $E = \frac{kQ}{r^{2}}$，式中 $k$ 为静电力常量，该均匀带电球所带的电荷量 $Q$ 为多大？
	\item
	求球心与球表面间的电势差 $\Delta U$；
	\item
	质量为 $m$，电荷量为 $q$ 的负电荷在球面处需具有多大的速度可以刚好运动到 $2R$ 处？
\end{enumerate}
\onepicture[align=right, w=4.5cm]{\includesvg{figs/32}}
%% answer: $\UV$（$\UNcdotm/\UC$）；$Q = \frac{E_{0}R^{2}}{k}$；$\Delta U = \frac{1}{2}E_{0}R$；$v = \sqrt{\frac{qE_{0}R}{m}}$
\jdanswer{
\begin{enumerate}
	\item
	$\UV$（$\UNcdotm/\UC$）

	\item
	$Q = \frac{E_{0}R^{2}}{k}$

	\item
	$\Delta U = \frac{1}{2}E_{0}R$

	\item
	$v = \sqrt{\frac{qE_{0}R}{m}}$

\end{enumerate}
}
%% memo:
\memoanswer{
（1）曲线下面积的单位为电场强度和长度单位的乘积，其单位为 $\UV$（$\UNcdotm/\UC$）。

（2）均匀带电球表面处场强 $E_{0}=\frac{kQ}{R^{2}}$，得出 $Q=\frac{E_{0}R^{2}}{k}$。

（3）球心与表面间电势差为将单位正电荷由球心移动到球表面过程中电场力的功，其大小即图中三角形面积，因此 $\Delta U=\frac{1}{2}E_{0}R$。

（4）由动能定理有 $-q\Delta U'=0-\frac{1}{2}mv^{2}$，由题意 $\Delta U'=\frac{1}{2}E_{0}R$，得 $v=\sqrt{\frac{qE_{0}R}{m}}$。
}

\item
%% number: 33
%% paperName: 2013年上海高考第 33 题 16 分
%% typeId: 6
%% type: 计算题
%% chapter: 电磁感应
%% point: 电磁感应能量分析
%% method: 无
%% score: 16
%% degree: 450
%% duplicateId: 0
%% body:
如图，两根相距 $l = 0.4\Um$、电阻不计的平行光滑金属导轨水平放置，一端与阻值 $R = 0.15\UO$ 的电阻相连。导轨 $x > 0$ 一侧存在沿 $x$ 方向均匀增大的稳恒磁场，其方向与导轨平面垂直，变化率 $k = 0.5\UTm$，$x = 0$ 处磁场的磁感应强度 $B_{0} = 0.5\UT$。一根质量 $m = 0.1\Ukg$、电阻 $r = 0.05\UO$ 的金属棒置于导轨上，并与导轨垂直。棒在外力作用下从 $x = 0$ 处以初速度 $v_{0} = 2\Ums$ 沿导轨向右运动，运动过程中电阻上消耗的功率不变。求：
\begin{enumerate}
	\item
	回路中的电流；
	\item
	金属棒在 $x = 2\Um$ 处的速度；
	\item
	金属棒从 $x = 0$ 运动到 $x = 2\Um$ 过程中安培力做功的大小；
	\item
	金属棒从 $x = 0$ 运动到 $x = 2\Um$ 过程中外力的平均功率。
\end{enumerate}
\onepicture[align=right, w=5.5cm]{\includesvg{figs/33}}
%% answer: $I=2\UA$；$v=\frac{2}{3}\Ums$；$W_{FA}=1.6\UJ$；$\overline{P}=0.7\UW$
\jdanswer{
\begin{enumerate}
	\item
	$I=2\UA$

	\item
	$v=\frac{2}{3}\Ums$

	\item
	$W_{FA}=1.6\UJ$

	\item
	$\overline{P}=0.7\UW$

\end{enumerate}
}
%% memo:
\memoanswer{
解：（1）电阻上消耗的功率恒定即电流恒定，在 $x=0$ 处有 $E=B_{0}lv_{0}=0.5\times0.4\times2\UV=0.4\UV$，

$I=\frac{E}{R+r}=\frac{0.4}{0.15+0.05}\UA=2\UA$

（2）由题意，磁感应强度 $B=B_{0}+kx$，考虑到电流恒定，有

$\frac{B_{0}lv_{0}}{R+r}=\frac{(B_{0}+kx)lv}{R+r}$

解得 $v=\frac{B_{0}v_{0}}{B_{0}+kx}=\frac{0.5\times2}{0.5+0.5\times2}\Ums=\frac{2}{3}\Ums$

（3）棒受到的安培力 $F_{A}=BIl=(B_{0}+kx)Il=0.4(1+x)\UN$，安培力随位置线性变化，

$W_{FA}=\frac{1}{2}\left[B_{0}+(B_{0}+kx)\right]Ilx$

代入数值后得 $W_{FA}=1.6\UJ$

（4）由动能定理 $W_{F}-W_{FA}=\frac{1}{2}mv^{2}-\frac{1}{2}mv_{0}^{2}$，其中外力做功

$W_{F}=W_{FA}+\frac{1}{2}mv^{2}-\frac{1}{2}mv_{0}^{2}=\left(1.6-\frac{16}{90}\right)\UJ=1.4\UJ$

安培力做功即电阻上消耗的能量，即 $W_{FA}=I^{2}(R+r)t$，运动时间为

$t=\frac{W_{FA}}{I^{2}(R+r)}=\frac{1.6}{2^{2}\times0.2}\Us=2\Us$

还可以使用以下方法求出运动时间：磁感应强度随位置线性变化，

$q=\frac{\Delta\Phi}{R+r}=\frac{(B_{0}+B_{0}+kx)lx}{2(R+r)}=\frac{(0.5+0.5+0.5\times2)\times0.4\times2}{2\times(0.15+0.05)}\UC=4\UC$

$t=\frac{q}{I}=\frac{4}{2}\Us=2\Us$

则 $\overline{P}=\frac{W_{F}}{t}=\frac{1.4}{2}\UW=0.7\UW$。
}

\end{enumerate}

\end{document}
